\clearpage\begingroup\SourceLeaf{656b}
\ArithmeticHead{21}{R}{II. DAS PRINZIP DES KLEINSTEN ZWANGES.}
\makeatletter\tagsleft@true\let\veqno\@@leqno\makeatother
\renewcommand{\thefootnote}{\arabic{footnote})}\setcounter{footnote}{37}
\begin{center}\textls[100]{2. Die \textsc{Gauss}sche Arbeit über das Prinzip des kleinsten Zwanges.}\end{center}
\noindent\hspace*{\parindent}\textsc{Gauss} geht nun einen Schritt weiter auf dem von \textsc{Lagrange} eingeschlagenen
Wege der Vereinheitlichung und Ökonomie, indem er den beiden letztgenannten
Prinzipien ein einziges neues, das des kleinsten Zwanges, substituiert. Er
empfindet es als Mangel, dass Statik und Dynamik gesondert behandelt werden,
während doch von einem höheren Standpunkte aus erstere nur einen ganz
speziellen Fall der Mechanik bilde. Diesen Mangel soll sein neues Prinzip
beheben, es soll also das einzige Gesetz der Mechanik schlechthin sein.

Es seien $A_1\, .\, .\, .\, A_n$ die Positionen der $n$ Massenpunkte mit den Koordinaten
$x_i$ im Zeitpunkte $t$, die einerseits den angreifenden Kräften mit den
Komponenten $X_i$ unterliegen, andererseits gewissen $m$ Bedingungsgleichungen
(1) genügen, deren Koeffizienten Funktionen der Koordinaten und der Zeit
sind. Die Bedingungsgleichungen äussern sich im Auftreten von Zwangskräften,
deren Zusammensetzung mit den angreifenden Kräften die wirklichen
Kräfte ergibt. Letzteren verdanken die Punkte ihre wirkliche Bewegung,
durch die sie im Zeitpunkte $t+\Delta t$ in die Stellungen $C_1\, .\, .\, .\, C_n$ gelangen, deren
Koordinaten durch $x_i+\dot x_i\Delta t+\frac12\ddot x_i(\Delta t)^2+ .\, .\, .$ gegeben sind; würden in den
Punkten $A_1\, .\, .\, .\, A_n$ nur die angreifenden Kräfte unter Ausschaltung der Zwangskräfte
wirken, so hätten die Punkte im Momente $t+\Delta t$ die Positionen $B_1\, .\, .\, .\, B_n$
mit den Koordinaten $x_i+\dot x_i\Delta t+\dfrac12\dfrac{X_i}{m_i}(\Delta t)^2+ .\, .\, .$ angenommen. Die Vektoren
$\overrightarrow{B_\nu C_\nu}\ (\nu=1\, .\, .\, .\, n)$ messen also durch ihre Länge die Abweichung der wirklichen
von der freien, d.\,h. ohne Bindungen vor sich gehenden Bewegung
nach Ablauf der infinitesimalen Zeit $\Delta t$. \textsc{Gauss} bezeichnet dann $m_{3\nu}|B_\nu C_\nu|^2$
als den Zwang, den der $\nu$-te Massenpunkt infolge der herrschenden Bedingungen
im Zeitelement $\Delta t$ erleidet, und somit
\[
Z=\sum_{\nu=1}^{n}m_{3\nu}|B_\nu C_\nu|^2
=\frac{(\Delta t)^4}{4}\sum_{i=1}^{3n}\frac1{m_i}(X_i-m_i\ddot x_i)^2\tag{5}
\]
\clearpage\SourceLeaf{657}\ArithmeticHead{22}{L}{H. GEPPERT, ÜBER GAUSS’ ARBEITEN ZUR MECHANIK UND POTENTIALTHEORIE.}
\noindent als den Zwang des ganzen Systems $\Sigma$. Die Vektoren $\overrightarrow{B_\nu C_\nu}$ sind andererseits
proportional zu Kräften, die sich vermöge der Bedingungen das Gleichgewicht
halten, d.\,h. infolge letzterer nicht wirksam werden. Wendet man also auf
sie das Prinzip der virtuellen Verrückungen an, so folgt, dass
\[
\sum_{i=1}^{3n}(X_i-m_i\ddot x_i)\delta x_i=0
\]
\noindent sein muss, womit sich das \textsc{d’Alembert}sche Prinzip ergibt.

\textsc{Gauss} gibt aber diesem letzten Schluss eine andere Wendung. Fasst
man nämlich $Z$ nur als Funktion der $\ddot x_i$ auf und sucht diejenigen Werte dieser
Variablen, die $Z$ zum Extremum (und zwar offenbar zum Minimum) machen,
so findet man die Bedingung
\[
\sum_{i=1}^{3n}(X_i-m_i\ddot x_i)\delta\ddot x_i=0,\tag{6}
\]
\noindent die formal mit der Gleichung des \textsc{d’Alembert}schen Prinzips (4) übereinstimmt,
wenn die $\delta\ddot x_i$ noch die den Beziehungen (3) entsprechenden Bedingungsgleichungen
\[
\sum_{i=1}^{3n}a_{\varkappa i}\delta\ddot x_i=0
\qquad\qquad(\varkappa=1\, .\, .\, .\, m)\tag{7}
\]
\noindent erfüllen. Diese letzten Gleichungen kann man aber folgendermassen deuten:
aus (1) folgt, dass die wirklichen Beschleunigungskomponenten $\ddot x_i$ $m$ Relationen
der Form
\[
\sum_{i=1}^{3n}a_{\varkappa i}\ddot x_i+\Phi_\varkappa(x_i,\dot x_i,t)=0
\qquad(\varkappa=1\, .\, .\, .\, m)\tag{8}
\]
\noindent genügen müssen, worin die Funktionen $\Phi_\varkappa$ die $\ddot x_i$ nicht enthalten. Ordnen
wir also der wirklichen Bahn eine variierte Bahn zu, bei der Koordinaten
und Geschwindigkeiten im Zeitpunkte $t$ unverändert bleiben:
\[
\delta x_i=\delta\dot x_i=0\qquad\qquad(i=1\, .\, .\, .\, 3n),
\]
\noindent und nur die Beschleunigungen um $\delta\ddot x_i$ variiert werden, und verlangen wir, dass
diese variierte Bahn ebenfalls zulässig sei, also den vorgelegten Bedingungen
(1) genüge, so muss neben (8) eben die Beziehung (7) für die Variationen
gelten. Man bezeichnet die dargelegte Art der Variation als eine \textsc{Gauss}sche
\clearpage\SourceLeaf{658}\ArithmeticHead{23}{R}{II. DAS PRINZIP DES KLEINSTEN ZWANGES.}
\noindent Variation. Somit bedeutet (6), dass bei einer \textsc{Gauss}schen Variation der Bahn
$\delta Z=0$ ist, d.\,h. dass die wirkliche Bahn im Vergleich zu den nach \textsc{Gauss}
variierten Bahnen den Zwang zu einem Minimum macht. Dies ist das Prinzip
des kleinsten Zwanges, und es ist aus seiner Ableitung evident, dass es mit
dem Prinzip von \textsc{d’Alembert} und dem der virtuellen Verrückungen — wenigstens
solange die obige Ableitung ihre Geltung behält — völlig gleichwertig ist.

\textsc{Mach}\footnote{a.\,a.\,O. S. 338.} hat darauf hingewiesen, dass dieses Prinzip mit folgender Aussage
identisch ist: die Arbeit, durch die in einem Zeitelement die Abweichung
der wirklichen, d.\,h. erzwungenen Bewegung von der freien hervorgebracht
wird, also die sog. verlorene Arbeit, ist ein Minimum. Man kann aus dem
\textsc{Gauss}schen Prinzip eine Reihe weiterer Minimumsprinzipe ableiten, die alle
mit ihm äquivalent sind\footnote{Die vollständige Literatur ist aufgezählt bei \textsc{Voss} a.\,a.\,O. S. 84.}. In der letztgenannten Fassung erscheint das \textsc{Gauss}sche
Prinzip zugleich in neuer historischer Beleuchtung, denn schon \textsc{Carnot}\footnote{L. N. M. \textsc{Carnot}, \textit{Principes fondamentaux de l’équilibre et des mouvements.} Paris 1803, S. 157.} hatte
beim Versuche einer sachlichen Erklärung der metaphysischen Gedankengänge,
mit denen \textsc{Maupertuis} sein Prinzip der kleinsten Aktion umrahmte, darauf
hingewiesen, dass beim elastischen und unelastischen Stoss die verlorenen
lebendigen Kräfte ein Minimum sein müssten; insofern erscheint das \textsc{Gauss}sche
Prinzip als der Abschluss einer mit dem \textsc{Maupertuis}schen Prinzip zusammenhängenden
Entwicklung.

\textsc{Gauss} selbst hat sein neues Prinzip zum ersten Male in dem oben zitierten
Briefe an \textsc{Olbers} 1829 ausgesprochen und im gleichen Jahre in seiner
Abhandlung über diesen Gegenstand bekannt gemacht. In seinen späteren
Arbeiten ist er nicht mehr darauf zurückgekommen, und dies mag es wohl
z.\,T. erklären, dass sein Grundsatz erst erheblich später in seiner fundamentalen
Bedeutung erkannt und angewandt wurde. Noch kurz vor seinem Tode
entstand unter \textsc{Gauss}’ unmittelbarem Einfluss die Dissertation von A. \textsc{Ritter},
die neben einer ausführlichen Erläuterung Beispiele entwickelt und die elementaren
Sätze der Statik aus dem \textsc{Gauss}schen Prinzip ableitet. Das erste
Lehrbuch, das das \textsc{Gauss}sche Prinzip unter die dynamischen Grundgesetze
einreihte, war das von \textsc{Earnshaw} (\textit{Dynamics}, Cambridge 1839). Die ersten
Arbeiten über das \textsc{Gauss}sche Prinzip befassen sich nur mit Erläuterungen
\clearpage\SourceLeaf{659}\ArithmeticHead{24}{L}{H. GEPPERT, ÜBER GAUSS’ ARBEITEN ZUR MECHANIK UND POTENTIALTHEORIE.}
\noindent und direkten Anwendungen, so z.\,B. gibt \textsc{Reuschle}\footnote{\textsc{Reuschle}, \textit{Über das Prinzip des kleinsten Zwanges und die damit zusammenhängenden mechanischen Prinzipe.} Archiv d. Math. und Phys., 6, 1845, S. 238—270.} die Anwendung auf
den Hebel und \textsc{Scheffler}\footnote{H. \textsc{Scheffler}, \textit{Über das Gausssche Grundgesetz der Mechanik.} Zeitschr. f. Math. und Phys.,
3, 1858, S. 197—223, 261—274.} ausserdem die Behandlung der Pendelbewegung,
des unelastischen Stosses und der Bewegung eines schweren Punktes auf
einer Fläche. Diesen Beispielen hat \textsc{Mach} einige weitere aus der Statik und
Dynamik hinzugefügt. Der erste, der das \textsc{Gauss}sche Prinzip in seiner vollen
Tragweite würdigte, war wohl \textsc{Heinrich Hertz}\footnote{H. \textsc{Hertz}, \textit{Die Prinzipien der Mechanik.} Leipzig 1894. Ges. Werke III, vergl. S. 184 ff.}, denn in seinem Buche ist
dieses Prinzip zusammen mit dem \textsc{Galilei}schen Trägheitsgesetz zum Grundpfeiler
seiner kräftelosen Mechanik erhoben.

Der von \textsc{Gauss} gegebene Beweis seines Prinzips, der sich nur in der
geometrischen Einkleidung von den oben wiedergegebenen Überlegungen unterscheidet,
impliziert eine Voraussetzung, die \textsc{Gauss} selbst nicht beachtet, und
auf die erst \textsc{Stäckel}\footnote{P. \textsc{Stäckel}, \textit{Bemerkungen zum Prinzip des kleinsten Zwanges.} Sitzungsber. d. Heidelb. Akad.
d. Wiss. Abt. A. Jahrg. 1919.} aufmerksam gemacht hat, nämlich, dass alle auftretenden
Bedingungsgleichungen (1) linear in den Differentialen seien. Betrachtet
man die $m$-reihige Matrix
\[
\left(\begin{array}{@{}cccc@{}}
a_{11},&a_{12},&.\, .\, .&a_{1,3n}\\
a_{21},&a_{22},&.\, .\, .&a_{2,3n}\\
\multicolumn{4}{c}{.\; .\; .\; .\; .\; .\; .\; .\; .\; .}\\
a_{m1},&a_{m2},&.\, .\, .&a_{m,3n}
\end{array}\right),
\]
\noindent die selbst Funktion des Ortes und der Zeit ist, so kann sie entweder den
Rang $m$, oder einen Rang $r<m$ haben. Ist das Erstere der Fall, so nennen wir
die Lage des Systems $\Sigma$ regulär, ist hingegen $r<m$, so heisst die Lage
singulär. Tritt letzteres ein, so hat dies zur Folge, dass eine der Gleichungen
(1) durch eine Gleichung höheren Grades in den Differentialen zu ersetzen
ist, wie man daraus erkennt, dass es sich in den Bedingungen (1) ja um die
ersten Glieder von Taylorentwicklungen handelt. Nimmt man z.\,B. einen
Punkt mit der Masse Eins, der auf dem Kegel
\[
x_1^2+x_2^2-x_3^2=0
\]
\clearpage\SourceLeaf{660}\ArithmeticHead{25}{R}{II. DAS PRINZIP DES KLEINSTEN ZWANGES.}
\HistoricalFoot{X 2 Abh. 7.}{4}
\noindent zu laufen gezwungen ist, so lautet die Bedingung (1) im allgemeinen
\[
x_1dx_1+x_2dx_2-x_3dx_3=0.
\]
\noindent Sie versagt jedoch in der Kegelspitze, dem Orte der singulären Lage, und
statt ihrer müssen die Differentiale dort die Beziehung erfüllen
\[
(dx_1)^2+(dx_2)^2-(dx_3)^2=0.
\]
\noindent Für reguläre Lagen sind nun das \textsc{d’Alembert}sche und das \textsc{Gauss}sche Prinzip
gleichwertig; beide bestimmen die Beschleunigungen eindeutig aus dem augenblicklichen
Bewegungszustand, und auf solche Lagen bezog sich der obige
Beweis. Anders verhält es sich für singuläre Lagen; die einfachsten Beispiele
zeigen, dass in solchen Fällen das \textsc{d’Alembert}sche Prinzip versagt und zu
Widersprüchen führen kann; die Minimumsaufgabe des \textsc{Gauss}schen Prinzips
behält dagegen auch dann ihren Sinn, wenn die Variationen der Beschleunigungskomponenten,
die von der wirklichen zu der nach \textsc{Gauss} variierten Bahn
führen, nicht durch lineare, sondern allgemeiner durch algebraische Gleichungen
verknüpft sind. Ein Unterschied muss jedoch hervorgehoben werden: bei regulärer
Lage liefert das \textsc{Gauss}sche Prinzip ein einziges absolutes Minimum des Zwanges,
legt also die Beschleunigungen vollkommen eindeutig fest; bei singulärer Lage
kann es hingegen mehrere Minima des Zwanges geben, die die von dieser
Lage aus möglichen Bewegungen charakterisieren; der sog. Eindeutigkeitssatz
der Mechanik wird also für singuläre Lagen ungültig. Dass dem so ist, erkennt
man leicht, wenn man die Grössen $\ddot x_i\sqrt{m_i}$ als Koordinaten eines $3n$-dimensionalen
linearen Raumes $R_{3n}$ deutet; ist dann die Lage regulär, führen also die Bedingungen
(1) zu linearen Bindungen zwischen den $\ddot x_i$, so wird durch diese
eine im $R_{3n}$ eingebettete lineare Mannigfaltigkeit $R_{3n-m}$ ausgesondert, und
das \textsc{Gauss}sche Prinzip erfordert, geometrisch gedeutet, die Bestimmung des
kürzesten Abstandes des Punktes mit den Koordinaten $\dfrac{X_i}{\sqrt{m_i}}$ von diesem $R_{3n-m}$;
ist die Lage singulär, bestehen also zwischen den $\ddot x_i$ nichtlineare algebraische
Bedingungsgleichungen, so sondern diese aus dem $R_{3n}$ eine algebraische nichtlineare
Mannigfaltigkeit $V_{3n-m}$ aus, und die Bestimmung des kürzesten Abstandes
derselben von dem Punkte $\dfrac{X_i}{\sqrt{m_i}}$ kann dann zu mehreren Lösungen führen\footnote{Die Frage nach der physikalischen Möglichkeit oder Unmöglichkeit singulärer Lagen liegt ausserhalb
des Rahmens unserer Aufgabe; sie hängt damit zusammen, dass den mechanischen Prinzipien eine
unter Umständen physikalisch unzulässige Idealisierung zugrunde liegt.}.
\clearpage\SourceLeaf{661}\ArithmeticHead{26}{L}{H. GEPPERT, ÜBER GAUSS’ ARBEITEN ZUR MECHANIK UND POTENTIALTHEORIE.}
\noindent Somit folgt, dass für singuläre Lagen das \textsc{Gauss}sche Prinzip mehr leistet als
das \textsc{d’Alembert}sche und also aus diesem nicht abgeleitet werden kann. Man
wird vielmehr das Prinzip des kleinsten Zwanges für singuläre Lagen als ein
neues, weitertragendes Axiom aufzufassen haben.
\clearpage\endgroup
